\section{Correctness and Parameter Sizes}

As opposed to the spilt trapdoor technique in \cite{AC:ALLW25} based on the more advanced module lattice assumption, the proposed protocol in this thesis underlied hardness assumption on integer lattice, which is a trivial case of aforementioned problem. 
% Put this before the table, only once
\newcounter{paramrow}
\newcommand{\rowidx}[1]{%
    \refstepcounter{paramrow}%
    \label{#1}%
    [\theparamrow]%
}
\newcommand{\paramref}[1]{parameter~[\ref{#1}]}

{\scriptsize
\setlength{\tabcolsep}{2pt}
\renewcommand{\arraystretch}{0.78}
\setcounter{paramrow}{0}

\begin{longtable}{@{}
>{\raggedright\arraybackslash}p{0.07\textwidth}
>{\raggedright\arraybackslash}p{0.52\textwidth}
>{\raggedright\arraybackslash}p{0.30\textwidth}
>{\raggedleft\arraybackslash}p{0.05\textwidth}
@{}}

\caption{Parameter settings for the proposed construction.}
\label{tab:parameters}
\\

$q$
&
\(\begin{gathered}
q\ \text{prime},\\[-0.2em]
n/q^{m-n+1}\le \mathrm{negl}(\lambda),\\[-0.2em]
nk/q^{2mk-nk+1}\le \mathrm{negl}(\lambda),\\[-0.2em]
nk/q^{\overline m-nk+1}\le \mathrm{negl}(\lambda)
\end{gathered}\)
&
Modulus and full-rank/primitive conditions
&
\rowidx{param:q-prime}
\\

&
\(\begin{gathered}
q
\ge
O\!\left(\ell^{3/2}t^{7/2}\right)
\cdot
\omega\!\left(\log(2\ell t+\widetilde m)\lambda\log(t)\right),
\\[-0.2em]
q>
4\!\left(
\chi_1+
\beta\chi_0\sqrt{2mk}
+
(t+\overline m)\sigma_1\chi_0'
\right)
\end{gathered}\)
&
Modulus size and correctness requirement
&
\rowidx{param:q-bound}
\\

$n$
&
\(\ge \lambda\)
&
Base matrix row parameter
&
\rowidx{param:n}
\\

\(m_{\mathbf{G}}\)
&
\(= n\cdot \lceil \log q\rceil\)
&
Width of the gadget matrix \(\mathbf{G}_n\)
&
\rowidx{param:mG}
\\

\(m\)
&
\(= m_{\mathbf{G}}+n\)
&
Base matrix column parameter
&
\rowidx{param:m}
\\

\(\widetilde m\)
&
\(= nk\cdot \lceil \log q\rceil\)
&
Width of the gadget matrix \(\mathbf{G}_{nk}\)
&
\rowidx{param:mtilde}
\\

\(\overline m\)
&
\(\ge 3nk\lceil\log q\rceil\)
&
Width of \(H_{\mathsf{ro}}(r)\in\mathbb{Z}_q^{nk\times\overline m}\)
&
\rowidx{param:mbar}
\\

\(\ell\)
&
\(= poly(\lambda)\)
&
Batch size bound
&
\rowidx{param:ell}
\\

\(d\)
&
\(= O(\log|\mathcal I|)\)
&
Identity-space parameter
&
\rowidx{param:d}
\\

\(t\)
&
\(= \widetilde m(d+1)\)
&
Shifted multi-preimage sampler dimension
&
\rowidx{param:t}
\\


\(\eta_{\mathbf C}\)
&
\(\ge 8\sqrt{2mk}\cdot q^{n/(2m)+1/(mk)}\)
&
Upper bound of \(\eta_\varepsilon(\Lambda_q^\perp(\mathbf C))\)
&
\rowidx{param:eta-C}
\\

\(\eta_{\widetilde{\mathbf C}}\)
&
\(\ge 8\sqrt m\cdot q^{n/m+2/m}\)
&
Upper bound of \(\eta_\varepsilon(\Lambda_q^\perp(\mathbf C_j))\)
&
\rowidx{param:eta-Ctilde}
\\

\(\eta_{\mathbf G}\)
&
\(\ge \sqrt5\cdot \omega(\sqrt{\log(m_{\mathbf G})})\)
&
Upper bound of \(\eta_\varepsilon(\Lambda_q^\perp(\mathbf G_n))\)
&
\rowidx{param:eta-G}
\\

\(B_{\mathbf C}\)
&
\(\ge \eta_{\mathbf C}\cdot
\sqrt{2m_{\mathbf G}(m-m_{\mathbf G})k}\)
&
Upper bound of \(\lVert\mathbf T_{\mathbf C}\rVert\)
&
\rowidx{param:B-C}
\\

\(B_{\widetilde{\mathbf C}}\)
&
\(\ge \eta_{\widetilde{\mathbf C}}\cdot
\sqrt{2m_{\mathbf G}(m-m_{\mathbf G})}\)
&
Upper bound of \(\lVert\mathbf T_{\mathbf C_j}\rVert\)
&
\rowidx{param:B-Ctilde}
\\

\(B_{\mathbf U}\)
&
\(\ge s_{\max}(\mathbf S_{\mathrm{Cor}})\cdot\sqrt{m^2k}\)
&
Upper bound of \(\lVert\mathbf U_j\rVert\)
&
\rowidx{param:B-U}
\\

\(\mathbf S\)
&
\(=\mathbf S_{\mathrm{Cor}}\)
&
Gaussian width of \(\mathsf{PTrapGen}\)
&
\rowidx{param:S}
\\

\(\widetilde{\mathbf S}_i\)
&
\(\widetilde{\mathbf S}_i=\mathbf S_{\mathrm{Cor}}
\quad \forall i\in[m(k-1)]\)
&
Assumption parameter in~\cite{AC:ALLW25}
&
\rowidx{param:S-tilde-cor}
\\

&
\(\widetilde{\mathbf S}_i=\mathbf S_{\mathrm{Hon}}
\quad
\forall i\in[2m(k-1)]\setminus[m(k-1)]\)
&
&
\rowidx{param:S-tilde-hon}
\\

\(\mathbf S_{\mathrm{Hon}}\)
&
\(\begin{gathered}
2m(k-1)\text{-}\mathrm{SIS}_{\mathbb Z_q,n,2mk,
\{\widetilde{\mathbf S}_i\}_{i},\alpha}
\text{ is hard}
\end{gathered}\)
&
Assumption parameter in~\cite{AC:ALLW25}
&
\rowidx{param:S-hon-SIS}
\\

&
\(\begin{gathered}
2m(k-1)\text{-}\mathrm{LWE}_{\mathbb Z_q,n,2mk,
\{\widetilde{\mathbf S}_i\}_{i},\chi_0}
\text{ is hard}
\end{gathered}\)
&
&
\rowidx{param:S-hon-LWE}
\\

&
\(\max\{\eta_{\mathbf C},2q^{n/(2m)}\}
< s_{\min}(\mathbf S_{\mathrm{Hon}})\)
&
&
\rowidx{param:S-hon-smin}
\\

&
\(s_{\max}(\mathbf S_{\mathrm{Hon}})
\le
\min\{q/\sqrt m,\;s_{\min}(\mathbf S_{\mathrm{Hon}})\}\)
&
&
\rowidx{param:S-hon-smax}
\\

\(\mathbf S_{\mathrm{Cor}}\)
&
\(\begin{gathered}
2m(k-1)\text{-}\mathrm{SIS}_{\mathbb Z_q,n,2mk,
\{\widetilde{\mathbf S}_i\}_{i},\alpha}
\text{ is hard}
\end{gathered}\)
&
Assumption parameter in~\cite{AC:ALLW25}
&
\rowidx{param:S-cor-SIS}
\\

&
\(\begin{gathered}
2m(k-1)\text{-}\mathrm{LWE}_{\mathbb Z_q,n,2mk,
\{\widetilde{\mathbf S}_i\}_{i},\chi_0}
\text{ is hard}
\end{gathered}\)
&
&
\rowidx{param:S-cor-LWE}
\\

&
\(s_{\min}(\mathbf S_{\mathrm{Cor}})
\ge 2B_{\mathbf C}\omega(\sqrt{\log(2mk)})\)
&
&
\rowidx{param:S-cor-smin-1}
\\

&
\(s_{\min}(\mathbf S_{\mathrm{Cor}})
\ge \eta_{\mathbf C}\)
&
&
\rowidx{param:S-cor-smin-2}
\\

&
\(\max\{\eta_{\mathbf C},2(a^{2m(k-1)}q^n)^{1/(2m)}\}
<
s_{\min}(\mathbf S_{\mathrm{Cor}})\)
&
&
\rowidx{param:S-cor-smin-3}
\\

&
\(s_{\max}(\mathbf S_{\mathrm{Cor}})
\le
\min\{q/\sqrt m,\;s_{\min}(\mathbf S_{\mathrm{Cor}})\}\)
&
&
\rowidx{param:S-cor-smax}
\\

&
\(\begin{aligned}[t]
s_{\min}(\mathbf S_{\mathrm{Cor}})
&\ge
4\eta_{\widetilde{\mathbf C}}
s_{\max}(\mathbf S_{\mathrm{Hon}})
m\sqrt{(m-m_{\mathbf G})k(k-1)}
\\[-0.2em]
&\quad{}\cdot
\omega\!\left(\sqrt{\log(2mk+m_{\mathbf G})}\right)
\end{aligned}\)
&
&
\rowidx{param:S-cor-basis-bound}
\\

\(\mathbf S_{\mathrm{BASIS}}\)
&
\(\displaystyle
=
\begin{pmatrix}
\mathbf S_{\mathrm{Cor}} & 0 \\
0 & \sigma_{\mathbf G}\mathbf I_{m_{\mathbf G}}
\end{pmatrix},
\quad
\sigma_{\mathbf G}\ge \eta_{\mathbf G}\)
&
Parameter for BASIS-style simulation in~\cite{AC:ALLW25}
&
\rowidx{param:S-basis}
\\

\(\sigma_1\)
&
\(\begin{aligned}[t]
\ge
\max\Big\{&
18(\ell t+\widetilde m+\ell\overline m)^{3/2}
\log\!\big((\ell t+\widetilde m+\ell\overline m)\lambda\big)
\log(\log q),
\\[-0.2em]
&
2\overline m^{3/2}
\cdot
\omega\!\left(\log(\overline m+t)\right),
\\[-0.2em]
&
\frac{q}{w}\log(2\ell t)
\Big\}
\end{aligned}\)
&
Gaussian parameter for explainable \(\mathsf{SampleLeft}\) in~\cite{AC:ALLW25}
&
\rowidx{param:sigma1}
\\

\(\sigma_2\)
&
\(\ge 2B_{\mathbf U}B_{\widetilde{\mathbf C}}
\omega(\sqrt{\log m})\)
&
Gaussian width of \(\mathsf{PSampPre}\)
&
\rowidx{param:sigma2-bound-1}
\\

&
\(\ge B_{\mathbf U}\eta_{\widetilde{\mathbf C}}\)
&
&
\rowidx{param:sigma2-bound-2}
\\

&
\(\ge 2B_{\mathbf U}\omega(\sqrt{\log(mk)})\)
&
&
\rowidx{param:sigma2-bound-3}
\\

\(\beta\)
&
\(\ge k\sigma_2\sqrt{2mk}\)
&
Correctness bound for the reconstructed \(\mathsf{sbk}\)
&
\rowidx{param:beta}
\\

\(\alpha\)
&
\(\begin{gathered}
2m(k-1)\text{-}\mathrm{SIS}_{\mathbb Z_q,n,2mk,
\{\widetilde{\mathbf S}_i\}_{i},\alpha}
\text{ is hard}
\end{gathered}\)
&
SIS hardness-assumption parameter
&
\rowidx{param:alpha}
\\

\(\chi_0\)
&
\(\begin{gathered}
2m(k-1)\text{-}\mathrm{LWE}_{\mathbb Z_q,n,2mk,
\{\widetilde{\mathbf S}_i\}_{i},\chi_0}
\text{ is hard}
\end{gathered}\)
&
LWE indistinguishability-assumption parameter
&
\rowidx{param:chi0}
\\

\(\chi_2\)
&
\(\ge 8\sqrt{2mk}\cdot q^{n/(2m)+1/(mk)}\)
&
Gaussian parameter used in the PT indistinguishability proof
&
\rowidx{param:chi2}
\\

\(\chi_1\)
&
\(\ge \lambda^{\omega(1)}\chi_0\chi_2\cdot 2mk\)
&
Noise-flooding parameter for the PT challenge
&
\rowidx{param:chi1}
\\

\(\chi_0'\)
&
\(>
O\!\left(
\lambda^{\omega(1)}
\chi_0mk\sqrt{t+\overline m}
\right)\)
&
Noise-flooding parameter for hiding \(\mathbf e_2^{\top}\mathbf H\)
&
\rowidx{param:chi0-prime}
\\

\(w\)
&
\(\begin{gathered}
\operatorname{rank}(\mathbf A_i)
=
\operatorname{rank}(H_{\mathsf{ro}}(r))
=
nk,
\\[-0.2em]
\lambda_1^\infty(\mathbf A_i),
\lambda_1^\infty(H_{\mathsf{ro}}(r))
\ge w
\end{gathered}\)
&
Full-rank and minimum-distance condition for all \(i\in[\ell]\), \(r\in\mathcal I\)
&
\rowidx{param:w}
\\


\end{longtable}
}
\section{Security}

We proof the security of scheme under the static corruption adversary with the highly-selective access to key query oracle. Let $L$ be the number of key querries that the adversary issues

\begin{figure}[h]
\centering
\scriptsize
\setlength{\fboxsep}{6pt}
\setlength{\fboxrule}{0.4pt}
\resizebox{\linewidth}{!}{%        <-- scales the box to the text width
\fbox{%
\begin{varwidth}{40cm}%           <-- large, so content sets its natural width
\begin{pchstack}
\procedure{{\normalsize$\mathsf{Exp}_{\Pi_{\mathsf{TBIBE}},\mathcal{A}}^{b}(1^\lambda,\ell,k,N)$}}{
    \rule{0pt}{0.6em}\\[-0.5em]
    (\mathcal{C},r^*) \gets \mathcal{A}(1^\lambda,1^\ell,1^k,1^N) \\
    \mathsf{assert}\ \mathcal{C}\subset_{<k}[N] \\
    (\mathsf{pk},(\mathsf{sk}_z)_{z\in[N]}) \gets \mathsf{TBIBE.Setup}(1^\lambda,1^\ell,1^k,1^N) \\
    \mathcal{Q} := \emptyset \\
    (T^*,z^*,m_0,m_1) \gets
        \mathcal{A}^{\mathsf{PreDecO}(\cdot)}(\mathsf{pk},(\mathsf{sk}_z)_{z\in\mathcal{C}}) \\
    (\mathsf{chal},\mathsf{ct}^*) \gets \mathsf{ChalO}(T^*,z^*,m_0,m_1) \\
    b' \gets \mathcal{A}^{\mathsf{PreDecO}(\cdot)}(\mathsf{chal},\mathsf{ct}^*) \\
    \pcreturn b'
}
\pchspace[1.2em]
\begin{pcvstack}
\procedure{{\normalsize$\mathsf{ChalO}(T^*,z^*,m_0,m_1)$}}{
    \rule{0pt}{0.6em}\\[-0.5em]
    \mathsf{assert}\ z^*\in T^*\setminus\mathcal{C} \\
    \pcfor z\in T^*\setminus\{z^*\}: \\
    \pcind \mathsf{sbk}^*_z \gets \mathsf{TBIBE.PreDec}(\mathsf{sk}_z,T^*,r^*) \\
    \mathsf{chal} := (\mathsf{sbk}^*_z)_{z\in T^*\setminus\{z^*\}} \\
    \mathsf{ct}^* \gets \mathsf{TBIBE.Enc}(\mathsf{pk},r^*,m_b) \\
    \pcreturn (\mathsf{chal},\mathsf{ct}^*)
        \hphantom{\mathsf{sbk}_{i,z}\gets \mathsf{TBIBE.PreDec}(\mathsf{sk}_z,T_i,\{r_{i,j}\}_{j\in[\ell]})}
}
\pcvspace[0.8em]
\procedure{{\normalsize$\mathsf{PreDecO}(T_i,\{r_{i,j}\}_{j\in[\ell]})$}}{
    \rule{0pt}{0.6em}\\[-0.5em]
    \mathsf{assert}\ T_i\subseteq_k [N] \\
    \pcif \{H(r_{i,j})\}_{j\in[\ell]}\ \text{not distinct} \pcthen \pcreturn \bot \\
    \pcif r^*\in\{r_{i,j}\}_{j\in[\ell]} \pcthen \pcreturn \bot \\
    \pcif \{r_{i,j}\}_{j\in[\ell]}\in\mathcal{Q} \pcthen \pcreturn \bot \\
    \mathcal{Q} := \mathcal{Q}\cup\{\{r_{i,j}\}_{j\in[\ell]}\} \\
    \pcfor z\in T_i: \\
    \pcind \mathsf{sbk}_{i,z}\gets \mathsf{TBIBE.PreDec}(\mathsf{sk}_z,T_i,\{r_{i,j}\}_{j\in[\ell]}) \\
    \pcreturn (\mathsf{sbk}_{i,z})_{z\in T_i}
}
\end{pcvstack}
\end{pchstack}%
\end{varwidth}%
}}
\caption{Selective security experiment for the TBIBE scheme.}
\label{fig:TBIBE-security-experiment}
\end{figure}

\begin{theorem}
    Suppose the parameters follow Table~\ref{tab:parameters}. Let $\mathcal{A}$ be a PPT adversary for the TBIBE security game (Figure~\ref{fig:TBIBE-security-experiment}), and $\mathcal{B}$ be the PPT adversaries for $(nk,2mk,\chi_0,\chi_1,\sigma_2)$$\epsilon$-indistinguishability security games (Theorem~\ref{theo:indistinguishibity_oneway}). THen we have:
    $$\mathsf{Adv}_{\mathcal{A}}^{\mathsf{TBIBE}}(1^\lambda,1^\ell,1^k,1^N) \le \mathsf{Adv}_{\mathcal{B}}^{\mathsf{ind}}(nk,2mk,\chi_0,\chi_1,\sigma_2) + \mathsf{negl}(\lambda) + \epsilon/2$$
    whenever $H_{\mathsf{sp}}, H_{\mathsf{ro}}$ are modeled as random oracles.
\end{theorem}
    

First, we represent a series of highly selective hybrid games and prove they are statistically equivalent up to a negligible difference. Assume that we have a PTT adversary $\mathcal{A}$ that:

\noindent\textbf{Hybrid 0:} This is the $\mathsf{Exp}_{\Pi_{\mathsf{TBIBE}},\mathcal{A}}^{b}$ in random oracle model
\begin{itemize}
    \item \textbf{Before the setup:}  The adversary $\mathcal{A}$ sends $(\mathcal{C}, r^*)$
to the PTT challenger where $\mathcal{C} \subset_{<k} [N]$ is the corruption set and $r^*$ is the challenge identity. The challenger then checks whether $\|\mathcal{C}\| < k$ is satisfied. If not, the PPT challenge halt the game and output random bit

\item \textbf{Setup:} The challenger caculates $\mathsf{crs}_{\mathsf{samp}} \leftarrow \mathsf{Gen}(1^\lambda,1^\ell)$. 
    It then samples a tuple 
    $(\mathbf{C},\mathbf{T}_{\mathbf{C}}) \leftarrow \mathsf{TrapGen}(1^{nk},q,2mk)$, 
    as well as a uniformly random vector $\mathbf{u} \in \mathbb{Z}_q^{nk}$. 
    It invokes $\mathsf{PTrapGen}(\mathbf{C}, \mathbf{T_C},j)$ to get $\mathsf{ptd}_j := (\mathbf{U}_j,\mathbf{C}_j, \mathbf{T}_{\mathbf{C}_j})$ for each $j \in T_i$ before sending 
    $\mathsf{pk}:=(\mathsf{crs}_{\mathsf{samp}}, \mathbf{C},\mathbf{u})$ to $\mathcal{A} \text{ and } (\mathsf{ptd}_j)_{j \in \mathcal{C}}$. Finally, it initializes an empty set $\mathcal{Q}$.

\item \textbf{Respond to pre-challenge queries:} For each $i \in [L]$, the challenger receives the incoming query $(T_i, \{r_{i,j}\}_{j\in[\ell]})$ in which $T_i$ is a set of active parties and $\{r_{i,j}\}_{j\in[\ell]}$ is an identity batch. Before proceeding, it first checks whether the inputs follow the below requirements:
\begin{itemize}
    \item All identities within the batch are collision-free under $H$
    \item The challenge identity $r^*$ is absent from the batch.
    \item The batch was not submitted in any prior query.
\end{itemize}

If all the above conditions pass, the challenger continues by adding this query to set $\mathcal{Q}$ and then calculating $H(r_{i,j})=\widetilde{h}_{i.j}$ for $j\in[\ell]$ and runs 
    $(\mathbf{A}_{i,1},\ldots,\mathbf{A}_{i,\ell},\mathsf{td_{i}}=(\mathbf{A}_i,\mathbf{T}_{\mathbf{A_i}}))
    =
    \mathsf{Expand}(1^\lambda,1^\ell,\mathsf{crs}_{\mathsf{samp}},
    (\widetilde{h}_i)_{i\in[\ell]})$ 
    to obtain the matrices $\mathbf{A}$ and $\widehat{\mathbf{A}}$. 
     To obtain $\mathsf{rb}$, it queries $H_{\mathsf{sp}}(\mathbf{A}_i,\mathbf{T}_{\mathbf{A}_i},(r_{i,j})_{j\in[\ell]},\mathbf{u},\sigma_1)$, and then invokes 
    $\mathsf{SampleLeft}(\mathbf{A},\widehat{\mathbf{A}},
    \mathbf{T}_{\mathbf{A}},\mathbf{h},\sigma_1;\mathsf{rb})$ 
    to obtain $\widehat{\mathbf{c}}$. After that, it returns 
    $\mathsf{sbk}_{z} \leftarrow 
    \mathsf{PSamplePre}(\mathbf{C}_z,\mathsf{ptd}_z,T_z,
    \mathbf{G}\widehat{\mathbf{c}},\sigma_2)$ for $z \in [N]$. In addition, for the challenge hint query, the challenger also does the same as above, except that it does not send the partial decryption key of the target party $z^*$.

\item \textbf{Challenge:} $\mathcal{A}$ sends the target active set $T*$, the party $z^*$ which belongs to $T^*$, along with two messages $m_0, m_1$ to the challenger and receives back the ciphertext tuples $\mathsf{ct}^* \leftarrow \mathsf{TBIBE.Enc}(\mathsf{pk}, r^*, m_b)$. 

\item \textbf{Post-challenge queries}: adversary $\mathcal{A}$ may again issue an arbitrary number of key computation queries as in the pre-challenge round.


\item \textbf{Random oracle queries:}
The challenger simulates the random oracle queries $H_{\mathsf{ro}}$ and $H_{\mathsf{sp}}$ consistently across all rounds as follows:
\begin{itemize}
    \item For $H_{\mathsf{ro}}$ queries: Upon receiving an identity $r$,challenger samples $\mathbf{H} \leftarrow \mathbb{Z}_q^{n \times \overline{m}}$, stores $H_{\mathsf{ro}} := H$, and then output $\mathbf{H}$
    \item For $H_{\mathsf{sp}}$ queries: Upon receiving a tuple $(\mathbf{A},\mathbf{T_A},(r_i)_{i \in [\ell]}, \mathbf{u}, \sigma_1)$, challenger samples $\mathbf{h} \leftarrow \{0,1\}^{\rho_{\mathsf{SL}}}$, stores $H_{\mathsf{rp}} := h$, and then outputs $\mathbf{h}$
\end{itemize}
\item \textbf{Output:} At any point of time, $\mathcal{A}$ can halt the game and output $b' \in \{0,1\}$.

\end{itemize} 
\medskip
\noindent\textbf{Hybrid 1:} The challenger samples uniformly random $[\mathbf{H_1}|\mathbf{H}_2]=\mathbf{H} \in \{0,1\}^{2mk \times (t + \overline{m})}$, where $\mathbf{H}_1\in \mathbb{Z}_q^{mk\times t},\mathbf{H}_2  \in \mathbb{Z}_q^{mk \times \overline{m}}$, and sets $\mathbf{A}_{r^*} = \mathbf{CH}_1$. After that, it calculates $\mathsf{crs_{samp}}$ by running  the algorithm $\mathsf{GenProg}(1^\lambda, 1^\ell, \tilde{h^*}, H(r^*),\mathbf{A}_{r^*})$.\\
\medskip
\noindent\textbf{Hybrid 2:} Challenger programs the random oracle $H_{\mathsf{ro}}$ so that $H_{\mathsf{ro}}(r^*):=\mathbf{CH_2}$ \\
\medskip
\noindent\textbf{Hybrid 3:} For every $H_{\mathsf{ro}}$ query with identity $r \neq r^*$, it programs $(H_{\mathsf{ro}},\mathbf{T}_r) \leftarrow \mathsf{TrapGen}(1^{nk},q,\overline{m})$\\
\medskip
\noindent\textbf{Hybrid 4:} When the adversary makes a fresh query
\((\mathbf{A},\mathbf{T}_{\mathbf{A}},\widehat{\mathbf{A}},\widehat{\mathbf{B}},(r_i)_{i\in[\ell]},\mathbf{u},\sigma_1)\) to the random oracle \(H_{\mathsf{sp}}\), the challenger first reconstructs the matrices
\(\widehat{\mathbf{A}}\) and \(\mathbf{A}\) using
\(\mathsf{crs}_{\mathsf{samp}}\), \(\widehat{\mathbf{A}}\), \(\widehat{\mathbf{B}}\), and the identity tags
\((r_i)_{i\in[\ell]}\). It then sets
\(\mathbf{F}:=[\mathbf{A}\mid\widehat{\mathbf{A}}]\) and
\(\mathbf{h}:=(\mathbf{u},\ldots,\mathbf{u})\in\mathbb{Z}_q^{\ell n}\), and samples
\(\mathbf{v}\leftarrow \mathbf{F}_{\sigma_1}^{-1}(\mathbf{h})\).
Next, the challenger computes
\(\mathsf{rb}\gets
\mathsf{ExplainSL}(1^\kappa,\mathbf{A},\widehat{\mathbf{A}},\mathbf{T}_{\mathbf{A}},
\mathbf{h},\mathbf{v},\sigma_1)\), with precision parameter $\kappa:=2q_\mathsf{sp}/\epsilon$.
Finally, it stores the tuple
\((\mathbf{A},\mathbf{T}_{\mathbf{A}},
(r_i)_{i\in[\ell]},\mathbf{u},\sigma_1,\mathbf{v},\mathsf{rb})\)
and returns \(\mathsf{rb}\) to \(\mathcal{A}\).
\noindent\begin{lemma}
Suppose the parameters follow Table~\ref{tab:parameters}. Then
\textbf{Hybrid 0} $\approx_s$ \textbf{Hybrid 1}.
\end{lemma}

\noindent\begin{proof}
The only difference between \textbf{Hybrid 0} and \textbf{Hybrid 1} is the
generation of the sampler CRS and the matrix associated with the challenge
identity $r^\ast$.

In \textbf{Hybrid 0}, the challenger samples
$
    \mathsf{crs}_{\mathsf{samp}}
    \gets
    \mathsf{Gen}(1^\lambda,1^\ell,1^d),
$
and the challenge matrix is obtained honestly as
\[
    \mathbf A_{r^\ast}
    :=
    \mathsf{ExpandLocal}
    (1^\lambda,\mathsf{crs}_{\mathsf{samp}},H(r^\ast)).
\]
In \textbf{Hybrid 1}, the challenger first samples
$
    \mathbf H_1 \gets \{0,1\}^{2mk\times t}
$
and sets
$
    \mathbf A_{r^\ast}:=\mathbf C\mathbf H_1
    \in \mathbb Z_q^{nk\times t}.
$
It then computes the programmed CRS
\[
    \mathsf{crs}_{\mathsf{samp}}^\ast
    :=
    \mathsf{GenProg}
    (1^\lambda,1^\ell,1^d,H(r^\ast),\mathbf A_{r^\ast}).
\]

We first show that the programmed challenge matrix
$\mathbf A_{r^\ast}=\mathbf C\mathbf H_1$ is statistically close to a
uniform matrix in $\mathbb Z_q^{nk\times t}$. By Table~\ref{tab:parameters},
we have
$
    \mathbf C\in \mathbb Z_q^{nk\times 2mk}.
$
Moreover, since
$
    m_{\mathbf G}=n\lceil\log q\rceil,
    m=m_{\mathbf G}+n,
$
we have
$
    2mk
    =
    2k(m_{\mathbf G}+n)
    \ge
    2nk\log q.
$
Thus the leftover hash lemma applies to the matrix
$\mathbf C\in\mathbb Z_q^{nk\times 2mk}$ and to each binary column of
$\mathbf H_1$. Therefore,
$
    (\mathbf C,\mathbf C\mathbf h)
    \approx_s
    (\mathbf C,\mathbf a),
$
where $\mathbf h\gets\{0,1\}^{2mk}$ and
$\mathbf a\gets\mathbb Z_q^{nk}$. Applying this argument independently to
the $t$ columns of $\mathbf H_1$, and using that $t=poly(\lambda)$ by
Table~\ref{tab:parameters}, gives
$
    (\mathbf C,\mathbf C\mathbf H_1)
    \approx_s
    (\mathbf C,\mathbf A),
$
where
$
    \mathbf A\gets\mathbb Z_q^{nk\times t}.
$

Next, by the somewhere programmability of $\Pi_{\mathsf{samp}}$ stated in
Lemma~\ref{lem:smps}, if the programmed matrix
$\mathbf A_{r^\ast}$ is uniform over $\mathbb Z_q^{nk\times t}$, then
$
    \mathsf{GenProg}
    (1^\lambda,1^\ell,1^d,H(r^\ast),\mathbf A_{r^\ast})
    \approx_s
    \mathsf{Gen}(1^\lambda,1^\ell,1^d).
$
Furthermore, the same lemma guarantees that the programmed opens
correctly at the challenge identity, namely
$
    \mathsf{ExpandLocal}
    (1^\lambda,\mathsf{crs}_{\mathsf{samp}}^\ast,H(r^\ast))
    =
    \mathbf A_{r^\ast}.
$

The table conditions used here are precisely:
\[
    q \text{ is prime},\qquad
    n\ge \lambda,\qquad
    m_{\mathbf G}=n\lceil\log q\rceil,\qquad
    m=m_{\mathbf G}+n,
\]
which imply
$
    2mk\ge 2nk\log q,
$
allowing the leftover hash lemma to be applied to
$\mathbf C\in\mathbb Z_q^{nk\times 2mk}$. In addition,
$
    t=\widetilde m(d+1)=poly(\lambda)
$
ensures that the union bound over the $t$ columns of
$\mathbf C\mathbf H_1$ preserves negligible statistical distance.

Combining the leftover-hash argument with the somewhere programmability of
$\Pi_{\mathsf{samp}}$, the joint view of the adversary in \textbf{Hybrid 0}
is statistically indistinguishable from its view in \textbf{Hybrid 1}.
Hence,
$
    \textbf{Hybrid 0}\approx_s \textbf{Hybrid 1}.
$
\end{proof}

\noindent\begin{lemma}
Suppose the parameters follow Table~\ref{tab:parameters}. Then, for all
PPT adversaries $\mathcal A$, \textbf{Hybrid 1} $\approx_s$ \textbf{Hybrid 2}.

\end{lemma}

\noindent\begin{proof}
The only difference between \textbf{Hybrid 1} and \textbf{Hybrid 2} is the
way the random-oracle value $H_{\mathsf{ro}}(r^\ast)$ is generated. In
\textbf{Hybrid 1}, it is sampled uniformly as
$
    H_{\mathsf{ro}}(r^\ast)
    \gets
    \mathbb Z_q^{nk\times \overline m},
$
whereas in \textbf{Hybrid 2}, the challenger samples
$
    \mathbf H_2 \gets \{0,1\}^{2mk\times \overline m}
$
and programs
$
    H_{\mathsf{ro}}(r^\ast)
    :=
    \mathbf C\mathbf H_2.
$

By the parameter choice in Table~\ref{tab:parameters}, the matrix
$\mathbf C\in \mathbb Z_q^{nk\times 2mk}$ has enough columns for
Lemma~\ref{lem:leftoverHash} to apply. Hence, for each column
$\mathbf h\gets\{0,1\}^{2mk}$ of $\mathbf H_2$, we have
$
    (\mathbf C,\mathbf C\mathbf h)
    \approx_s
    (\mathbf C,\mathbf c),
$
where $\mathbf c\gets\mathbb Z_q^{nk}$ is uniform.

Applying this argument to all $\overline m=poly(\lambda)$ columns of
$\mathbf H_2$ and taking a union bound, we obtain
$
    (\mathbf C,\mathbf C\mathbf H_2)
    \approx_s
    (\mathbf C,\mathbf C'),
$
where
$
    \mathbf C'\gets\mathbb Z_q^{nk\times \overline m}.
$
Therefore, the programmed value $\mathbf C\mathbf H_2$ in
\textbf{Hybrid 2} is statistically indistinguishable from the uniform
value of $H_{\mathsf{ro}}(r^\ast)$ in \textbf{Hybrid 1}. All other parts
of the adversary's view are generated identically, so
$
    \textbf{Hybrid 1}\approx_s \textbf{Hybrid 2}.
$
\end{proof}
\noindent\begin{lemma}
Suppose that $n\ge \lambda$ and
$\overline{m}\ge O(nk\log q)$. Then, for all PPT adversaries
$\mathcal A$,
$
    \textbf{Hybrid 2} \approx_s \textbf{Hybrid 3}.
$
\end{lemma}

\noindent\begin{proof}
The only difference between \textbf{Hybrid 2} and \textbf{Hybrid 3} is how
the random-oracle values $H_{\mathsf{ro}}(r)$ are generated for identities
$r\neq r^\ast$. In \textbf{Hybrid 2}, these values are sampled uniformly as
$
    H_{\mathsf{ro}}(r)
    \gets
    \mathbb Z_q^{nk\times \overline m},
$
whereas in \textbf{Hybrid 3}, they are generated together with gadget
trapdoors:
$
    (H_{\mathsf{ro}}(r),\mathbf T_r)
    \gets
    \mathsf{TrapGen}(1^\lambda,1^{nk},1^{\overline m}).
$

By the parameter choice in Table~\ref{tab:parameters}, we have
$
    \overline m \ge O(nk\log q),
$
which is sufficient for the gadget-trapdoor generation lemma
(Lemma~\ref{lem:gadget_trapdoor}) to apply with row dimension $nk$ and
column dimension $\overline m$. Hence, the matrix output by
$\mathsf{TrapGen}$ is statistically close to uniform:
$
    H_{\mathsf{ro}}(r)
    \approx_s
    \mathbf H,
    \mathbf H\gets \mathbb Z_q^{nk\times \overline m}.
$

Therefore, for each identity $r\neq r^\ast$, the distribution of
$H_{\mathsf{ro}}(r)$ in \textbf{Hybrid 3} is statistically indistinguishable
from its distribution in \textbf{Hybrid 2}. The trapdoor $\mathbf T_r$ is
only kept internally by the challenger and is not given to the adversary.

Since $\mathcal A$ makes only polynomially many random-oracle queries, a
union bound over all such queries preserves negligible statistical distance.
All other parts of the adversary's view are identical in the two hybrids.
Thus,
$
    \textbf{Hybrid 2} \approx_s \textbf{Hybrid 3}.
$
\end{proof}
\noindent\begin{lemma}
Suppose the parameters follow Table~\ref{tab:parameters}. In particular,
assume that
$
    q>2,
    \rho
    =
    \operatorname{poly}
    (\lambda,nk,\ell t+\widetilde m,\ell\overline m,\log q),
$
and for every matrix $\mathbf A$ appearing in a pre-decryption query,
$
    \sigma_1
    \ge
    \|\mathbf T_{\mathbf A}\|_{\infty}
    \cdot
    \sigma_{\mathsf{loss}}.
$
Here, since the trapdoor $\mathbf T_{\mathbf A}$ generated by the sampler has
$\|\mathbf T_{\mathbf A}\|_{\infty}=1$, it suffices that
$
    \sigma_1
    \ge
    \sigma_{\mathsf{loss}},
$
where
$
    \sigma_{\mathsf{loss}}
    =
    18(\ell t+\widetilde m+\ell\overline m)^{3/2}
    \log\!\bigl((\ell t+\widetilde m+\ell\overline m)\lambda\bigr)
    \log(\log q).
$
Then, for all PPT adversaries $\mathcal A$,
$
    \Delta(\textbf{Hybrid 4},\textbf{Hybrid 5})
    \le
    \operatorname{negl}(\lambda)+\epsilon/2.
$
\end{lemma}

\noindent\begin{proof}
The only difference between \textbf{Hybrid 4} and \textbf{Hybrid 5} is how
the pair $(\mathbf v,\mathbf{rb})$ associated with an
$H_{\mathsf{sp}}$ query is generated.

Let
$
    \Pi_{\mathsf{SL}}
    =
    (\mathsf{SampleLeft},\mathsf{ExplainSL})
$
be the explainable sampler from Theorem~\ref{the:explainable}, instantiated
with Gaussian width $\sigma_1$. For a pre-decryption query, define
$
    \mathbf F := [\mathbf A\mid \widehat{\mathbf A}],
$
where
$
    \mathbf A\in
    \mathbb Z_q^{\ell nk\times(\ell t+\widetilde m)}
   \text{ and }
    \widehat{\mathbf A}\in
    \mathbb Z_q^{\ell nk\times \ell\overline m}.
$
Hence,
$
    \mathbf F
    \in
    \mathbb Z_q^{
        \ell nk\times(\ell t+\widetilde m+\ell\overline m)
    }.
$
By the choice of $\sigma_1$ in Table~\ref{tab:parameters}, the explainable
sampler is applicable to this matrix width.

In \textbf{Hybrid 4}, the challenger first samples
$
    \mathbf v
    \leftarrow
    \mathbf F^{-1}_{\sigma_1}(\mathbf h),
$
and then computes
$
    \mathbf{rb}'
    :=
    \mathsf{ExplainSL}
    \bigl(
        1^\kappa,
        \mathbf A,
        \widehat{\mathbf A},
        \mathbf T_{\mathbf A},
        \mathbf h,
        \mathbf v,
        \sigma_1
    \bigr).
$
Thus the relevant distribution is
\[
    \mathcal D_0
    :=
    \left\{
        (\mathbf v,\mathbf{rb}') :
        \mathbf v\leftarrow \mathbf F^{-1}_{\sigma_1}(\mathbf h),
        \;
        \mathbf{rb}'
        :=
        \mathsf{ExplainSL}
        (1^\kappa,\mathbf A,\widehat{\mathbf A},
        \mathbf T_{\mathbf A},\mathbf h,\mathbf v,\sigma_1)
    \right\}.
\]

By correctness of the explainable sampler, sampling
$
    \mathbf v
    \leftarrow
    \mathbf F^{-1}_{\sigma_1}(\mathbf h)
$
is statistically close to sampling uniform coins
$\mathbf{rb}\leftarrow\{0,1\}^{\rho}$ and outputting
$
    \mathbf v
    :=
    \mathsf{SampleLeft}
    (
        \mathbf A,
        \widehat{\mathbf A},
        \mathbf T_{\mathbf A},
        \mathbf h,
        \sigma_1;
        \mathbf{rb}
    ).
$
Therefore,
$
    \mathcal D_0 \approx_s \mathcal D_1,
$
where
\[
\begin{aligned}
    \mathcal D_1
    :=
    \left\{
        (\mathbf v,\mathbf{rb}') :
        \mathbf{rb}\leftarrow\{0,1\}^{\rho},
        \;
        \mathbf v :=
        \mathsf{SampleLeft}
        (
            \mathbf A,
            \widehat{\mathbf A},
            \mathbf T_{\mathbf A},
            \mathbf h,
            \sigma_1;
            \mathbf{rb}
        ),
        \right.\\
        \left.
        \mathbf{rb}'
        :=
        \mathsf{ExplainSL}
        (
            1^\kappa,
            \mathbf A,
            \widehat{\mathbf A},
            \mathbf T_{\mathbf A},
            \mathbf h,
            \mathbf v,
            \sigma_1
        )
    \right\}.
\end{aligned}
\]

Next, by the explainability property of $\Pi_{\mathsf{SL}}$, the randomness
recovered by $\mathsf{ExplainSL}$ is statistically close to fresh uniform
coins. Hence,
\[
    \Delta(\mathcal D_1,\mathcal D_2)
    \le
    \operatorname{negl}(\lambda)+\frac{1}{\kappa},
\]
where
$
    \mathcal D_2
    :=
    \left\{
        (\mathbf v,\mathbf{rb}') :
        \mathbf{rb}'\leftarrow\{0,1\}^{\rho},
        \;
        \mathbf v :=
        \mathsf{SampleLeft}
        (
            \mathbf A,
            \widehat{\mathbf A},
            \mathbf T_{\mathbf A},
            \mathbf h,
            \sigma_1;
            \mathbf{rb}'
        )
    \right\}.
$\\
This is exactly the distribution used in \textbf{Hybrid 5}, where
$H_{\mathsf{sp}}$ returns fresh random coins.

Finally, set
$
    \kappa := \frac{2q_{\mathsf{sp}}}{\epsilon}.
$, where $q_{\mathsf{sp}}$ is the number of $H_{\mathsf{sp}}$ queries.
Since $q_{\mathsf{sp}}=\operatorname{poly}(\lambda)$, this choice keeps
$\kappa=\operatorname{poly}(\lambda)$. Taking a union bound over all
$q_{\mathsf{sp}}$ random-oracle queries to $H_{\mathsf{sp}}$, we obtain
\[
\begin{aligned}
    \Delta(\textbf{Hybrid 4},\textbf{Hybrid 5})
    &\le
    q_{\mathsf{sp}}
    \left(
        \operatorname{negl}(\lambda)+\frac{1}{\kappa}
    \right)  \\
    &=
    \operatorname{negl}(\lambda)+\frac{\epsilon}{2}.
\end{aligned}
\]
This proves the claim.
\end{proof}
Now, we are ready to construct a threshold adversary $\mathcal{B}$ wrapping around the adversary $\mathcal{A}$. For the sake of contradiction, assume that we have an y $\mathcal{A}$ that can distinguish between the real highly selective experiment and the one where the ciphertexts are purely uniform random vectors. We construct a PPT $\mathcal{B}^\mathcal{A}$ against the $(\chi_0,\chi_1)$. Let  $\mathcal{B}^\mathcal{A}$ proceed as follows:


\noindent\textbf{Hybrid 5:}
\begin{itemize}
    \item $\text{Receive from } \mathcal{A} \text{ a set } C \subset [N] \text{ of corrupt users, a set }
\{(T_i,(r_{i,j})_{j \in \ell})\}_{i \in [L]}$ of all secret-key queries, and a challenge $(T^\ast, z^\ast, r^\ast).$
$\text{W.l.o.g., assume } |C| = k - 1.$
    \item Pass $C$ to the PT indistinguishability challenger, and receive $(\mathsf{crs_{samp}}, \mathbf{\tilde{A}},\mathbf{\tilde{B}},\mathbf{C},\mathbf{u})$ and partial trapdoors 
    $(\mathbf{U}_j)_{j \in C}$. Let $\mathsf{mpk} := \mathbf{A}$ and 
    $\mathsf{msk}_j := \mathsf{ptd}_j$ for all $j \in C$.


\item $\mathcal{B}$ pass $(T^*, z^*)$ to PT challenger, and receive $\left((\mathbf{x}_j^*)_{j\in T^*\setminus\{z^*\}}, \mathbf{y}^*, \mathbf{ch}_0, ch_1\right)$  and then set $\mathsf{sbk}^*_j := \mathbf{x}^*_j, \mathbf{u} := \mathbf{y}^*$. It then obtain challenge ciphertext by setting $(\mathsf{ct}[0],\mathsf{ct}[1],\mathsf{ct}[2]) = (ch_1 + m_b\cdot\lfloor q/2\rfloor,\mathbf{ch}_0, \mathbf{ch}_0\cdot \mathbf{H} + \mathbf{e}')$, where $\mathbf{e}'\leftarrow\mathcal{D}_{\mathbb{Z},\chi'}^{t + \overline{m}}$

\item $\mathcal{B}$ first checks $r^* \in (r_{i,j})_{j \in [\ell]}$ and makes decision as following cases:
\item \textbf{Case 1}: Suppose $r^* = r_{i,j^*}$ for some index $j^*$. 
\begin{itemize}
    \item $\mathcal{B}$ samples
$\mathbf{v}_{i,j^*}=\operatorname{col}(\mathbf{v}_{i,j^*}^1,\mathbf{v}_{i,j^*}^2)
\leftarrow D_{\mathbb{Z}_q^{2t},\sigma_1}$ and sets
$\mathbf{c}_i:=-\mathbf{F}_{r^*}\mathbf{v}_{i,j^*}+\mathbf{u}$, where
$\mathbf{F}_{r^*}=[\mathbf{A}_{r^*}| H_{\mathsf{ro}}(r^*)]$. It then samples $\widehat{\mathbf{c}_i}\leftarrow \mathbf{G}_{\sigma_1}^{-1}(\mathbf{c}_i)$
using the algorithm in Lemma~\ref{lem:gaget_matrix_sam} and sets $\mathbf{sbk}_i := \bot$.

\item For the remaining $j \in [\ell], j \neq j^*$, which means $r_{i,j} \neq r^*$, it finds $\mathbf{v}_{i,j}$ by calling
$\mathsf{SampleLeft}\!\left(
H_{\mathsf{ro}}(r_{i,j}),
\mathbf{A}_{r_{i,j}},
\mathbf{H}_2,
\mathbf{T}_{r_{i,j}},
\mathbf{u}-\mathbf{c}_i,
\sigma_1
\right)$ and then permutes the output components to obtain $\mathbf{v}_{i,j}:= \operatorname{col}(\mathbf{v}_{i,j}^1,\mathbf{v}_{i,j}^2)$
\end{itemize}
\item \textbf{Case 2}: $r^* \notin (r_{i,j})_{j \in [\ell]}$
\begin{itemize}
    \item Simulator $\mathcal{B}$ queries $\mathsf{PSampPreO}(T_i)$ to obtain 
    $((\mathbf{x}_{i,j})_{j \in T_i}, \mathbf{y}_i)$, where $\mathbf{y}_i$ is uniform. Let $\mathsf{sbk}_{i,j} := \mathbf{x}_{i,j}$, $\mathbf{c}_{i} := \mathbf{y}_{i}$, challenger  derives $\mathsf{sbk}_i$ by calling $\mathsf{TBIBE.Rec}((\mathsf{sbk_{i,j}})_{j \in T_i})$. It then samples $\mathbf{\hat{c}_i} \leftarrow \mathbf{G}^{-1}_{\sigma_1}(\mathbf{c}_i)$ using the algorithm. 
    \item For all $j \in [\ell]$, it invokes $\mathsf{SampleLeft}\!\left(
H_{\mathsf{ro}}(r_{i,j}),
\mathbf{A}_{r_{i,j}},
\mathbf{H}_2,
\mathbf{T}_{r_{i,j}},
\mathbf{u}-\mathbf{c}_i,
\sigma_1
\right)$ and then permutes the output to obtain $\mathbf{v}_{i,j}:= \operatorname{col}(\mathbf{v}_{i,j}^1,\mathbf{v}_{i,j}^2)$.
\end{itemize}




    
    
\item After that, the challenger sets
\(\mathbf{v}_i :=
\mathsf{col}(\mathbf{v}_{i,1}^{1},\ldots,\mathbf{v}_{i,\ell}^{1},
\widehat{\mathbf{c}},
\mathbf{v}_{i,1}^{2},\ldots,\mathbf{v}_{i,\ell}^{2})\).
It then obtains
\(\mathsf{rb}_i :=
\mathsf{ExplainSL}(1^\kappa,
\mathbf{A}_i,
\widehat{\mathbf{A}}_i,
\mathbf{T}_{\mathbf{A}_i},
\mathbf{h}=(\mathbf{u},\ldots,\mathbf{u}),
\mathbf{v}_i,
\sigma_1)\).
Finally, the challenger stores
\(\mathsf{st}_i :=
(\mathbf{A}_i,
\mathbf{T}_{\mathbf{A}_i},
((r_j)_{j\in[\ell]})_i,
T_i,
\mathbf{u},
\sigma_1,
\mathbf{v}_i,
\mathbf{sbk}_i,
\mathsf{rb}_i)\).

\item Upon any $\mathsf{rb}$ query from $\mathcal{A}$, if $\mathsf{rb} = \bot$, simulator samples a random $\mathbf{y} = \{0,1\}^\rho$ and set $\mathsf{rb} = y$. It return $\mathsf{rb}$ to $\mathcal{A}$
\end{itemize}

\noindent\textbf{Hybrid 6:} Challenger replace $c_1 \in \mathbb{Z}_q$

\noindent\begin{lemma}
Suppose the parameters follow Table~\ref{tab:parameters}. In particular,
assume that:
\begin{enumerate}
    \item By \paramref{param:q-prime}, $q>2$; by
    \paramref{param:t} and \paramref{param:mbar}, the dimensions
    $t$ and $\overline m$ are polynomially bounded; by
    \paramref{param:sigma1}, $\sigma_1\ge \omega(\sqrt{\log(nk)})$;
    and by \paramref{param:sigma2-bound-3},
    $\sigma_2 \ge 2B_{\mathbf U}\cdot \omega(\sqrt{\log(mk)})$.

    \item By \paramref{param:sigma1}, for all identity tags $r$ appearing
    in pre-decryption queries,
    $\sigma_1
    \ge
    2\overline m^{3/2}\|\mathbf T_r\|_{\infty}
    \cdot \omega(\log(\overline m+t))$,
    where $\mathbf T_r$ is the gadget trapdoor for
    $H_{\mathsf{ro}}(r)$ generated by $\mathsf{TrapGen}$.

    \item By \paramref{param:w}, the matrices $\mathbf A_{r_i}$ and
    $H_{\mathsf{ro}}(r_i)$ are full-rank of rank $nk$,
    $\lambda_1^\infty(\mathbf A_{r_i})$, and $
    \lambda_1^\infty(H_{\mathsf{ro}}(r_i))\ge w$.
    Moreover, by \paramref{param:sigma1},
    $\sigma_1\ge (q/w)\log(2\ell t)$.

    \item By \paramref{param:chi0-prime}, the noise-flooding parameter
    satisfies
    $\chi_0'
    >
    O(\lambda^{\omega(1)}\chi_0mk\sqrt{t+\overline m})$.
\end{enumerate}
Then, for all PPT adversaries $\mathcal A$,
\textbf{Hybrid 4} $\approx_s$ \textbf{Hybrid 5}.
\end{lemma}

\begin{proof}
The transition from \textbf{Hybrid 4} to \textbf{Hybrid 5} consists of
statistically simulating the challenge ciphertext component $\mathsf{ct}[3]$
and the corresponding $H_{\mathsf{sp}}$ responses.

First, consider the challenge ciphertext. In \textbf{Hybrid 4}, the simulator
constructs
$\mathsf{ct}[3]
:=
\mathbf{ch}_0^{\top}\mathbf H+\mathbf e'
=
\mathbf s^{\top}\mathbf C\mathbf H
+
\mathbf e_2^{\top}\mathbf H
+
\mathbf e'$,
where
$\mathbf e'\leftarrow D_{\mathbb Z,\chi_0'}^{t+\overline m}$.
Note that, we have
$\mathbf H\in\{0,1\}^{2mk\times(t+\overline m)}$ and 
$\mathbf e_2\leftarrow D_{\mathbb Z,\chi_0}^{2mk}$.
Therefore, using $\|\mathbf H\|\le \|\mathbf H\|_F$, we get
$\|\mathbf e_2^{\top}\mathbf H\|
\le
\|\mathbf e_2\|\cdot\|\mathbf H\|
\le
O(\chi_0\sqrt{2mk})\cdot\sqrt{2mk(t+\overline m)}
=
O(\chi_0mk\sqrt{t+\overline m})$.
By \paramref{param:chi0-prime} and Lemma~\ref{lem:̉noise_flooding},
$\mathbf e_2^{\top}\mathbf H+\mathbf e'
\approx_s
D_{\mathbb Z,\chi_0'}^{t+\overline m}$.
Hence, the distribution of $\mathsf{ct}[3]$ in \textbf{Hybrid 4} is
statistically close to the corresponding distribution in \textbf{Hybrid 5}.

It remains to argue that the simulator can answer the relevant
$H_{\mathsf{sp}}$ and pre-decryption queries consistently. Consider a
pre-decryption query containing identities $(r_i)_{i\in[\ell]}$.

\noindent\textbf{Case 1: $r^\ast=r_{i^\ast}$ for some $i^\ast$.}
We have
$\mathbf H=[\mathbf H_1\mid \mathbf H_2]$ with
$\mathbf H_1\in\{0,1\}^{2mk\times t}$ and
$\mathbf H_2\in\{0,1\}^{2mk\times\overline m}$.
Since $H_{\mathsf{ro}}(r^\ast)=\mathbf C\mathbf H_2$ and
$\mathbf A_{r^\ast}=\mathbf C\mathbf H_1$, we have
$\mathbf F_{r^\ast}
:=
[\mathbf A_{r^\ast}\mid H_{\mathsf{ro}}(r^\ast)]
=
[\mathbf C\mathbf H_1\mid \mathbf C\mathbf H_2]$.
The simulator samples $\mathbf v_{i^\ast}$ and defines
$\mathbf c:=-\mathbf F_{r^\ast}\mathbf v_{i^\ast}+\mathbf u$.
By the full-rank and minimum-distance condition in \paramref{param:w},
together with the lower bound
$\sigma_1\ge(q/w)\log(2\ell t)$ from \paramref{param:sigma1},
the value $\mathbf c$ is statistically close to uniform over
$\mathbb Z_q^{nk}$. Therefore, by Lemma~\ref{lem:gaget_matrix_sam},
$\widehat{\mathbf c}\leftarrow \mathbf G_{\sigma_1}^{-1}(\mathbf c)$
is statistically close to
$D_{\mathbb Z^{\widetilde m},\sigma_1}$.

\noindent\textbf{Case 2: $r^\ast\notin (r_i)_{i\in[\ell]}$.}
In this case, the simulator obtains a uniformly distributed target
$\mathbf c$ from the simulated partial-preimage oracle. Again, by
Lemma~\ref{lem:gaget_matrix_sam},
$\widehat{\mathbf c}\leftarrow \mathbf G_{\sigma_1}^{-1}(\mathbf c)$
is statistically close to
$D_{\mathbb Z^{\widetilde m},\sigma_1}$.

In both cases, for every identity $r_i\neq r^\ast$, the simulator uses the
trapdoor $\mathbf T_{r_i}$ of $H_{\mathsf{ro}}(r_i)$ and runs
$\mathsf{SampleLeft}
(H_{\mathsf{ro}}(r_i),\mathbf A_{r_i},\mathbf T_{r_i},
\mathbf u-\mathbf c,\sigma_1)$.
By \paramref{param:sigma1}, the Gaussian parameter $\sigma_1$ is large
enough for gadget-trapdoor-based $\mathsf{SampleLeft}$. Hence, after the
corresponding coordinate permutation, the output is statistically close to
$\mathbf F_{r_i,\sigma_1}^{-1}(\mathbf u-\mathbf c)$, where
$\mathbf F_{r_i}:=[\mathbf A_{r_i}\mid H_{\mathsf{ro}}(r_i)]$.
The coordinate permutation does not change the discrete Gaussian distribution,
since permutation matrices are orthogonal.

Therefore, the vector
$\operatorname{col}
(\mathbf v_1^1,\ldots,\mathbf v_\ell^1,
\widehat{\mathbf c},
\mathbf v_1^2,\ldots,\mathbf v_\ell^2)$
is statistically close to a sample from
$\mathbf F_{\sigma_1}^{-1}(\mathbf h)$, where
$\mathbf F:=[\mathbf A\mid\widehat{\mathbf A}]$ and
$\mathbf h:=(\mathbf u,\ldots,\mathbf u)\in\mathbb Z_q^{\ell nk}$.
Consequently, the value
$\mathsf{rb}:=
\mathsf{ExplainSL}
(1^\kappa,\mathbf A,\widehat{\mathbf A},
\mathbf T_{\mathbf A},\mathbf h,\mathbf v,\sigma_1)$
has the same distribution as in the real $H_{\mathsf{sp}}$ simulation, up to
negligible statistical distance.

Finally, the simulated partial-preimage oracle samples a uniformly random
target and generates partial preimages using $\mathsf{PSampPre}$ exactly as
in the corresponding hybrid. Hence, the oracle responses in \textbf{Hybrid 5}
are statistically indistinguishable from those in \textbf{Hybrid 4}. Since
$\mathcal A$ makes only polynomially many random-oracle and pre-decryption
queries, a union bound over all such queries preserves negligible statistical
distance. Therefore, \textbf{Hybrid 4} $\approx_s$ \textbf{Hybrid 5}.
\end{proof}

\noindent\begin{lemma}
    Suppose the parameter follows the table~\ref{tab:parameters} and the lattice trapdoor scheme $\mathsf{PT}$ has $(nk,2mk,\chi_0,\chi_1, \sigma_2)$-indistinguishability, the \textbf{Hybrid 5} and \textbf{Hybrid 6} are computationally close
\end{lemma}
\begin{proof}
    If $c_1$ from $\mathsf{PT}$ challenger equals $\mathbf{u^\top s} + e_1$, the $\mathcal{B}^\mathcal{A}$ perfectly simulates $c_1$ in \textbf{Hybrid 5}. Else, $c_1$ is uniformly random, thereby $\mathcal{B}^\mathcal{A}$ perfectly simulates $c_1$ in \textbf{Hybrid 6}. Therefore, the success probability of $\mathcal{B}^\mathcal{A}$ is the same as $\mathcal{A}$, which yields the desirable result.
\end{proof}
